\section{The Minimax Criterion under Affine Restrictions}
\label{sec:restrictions}

Compact affine restrictions $\mathcal{R} = \{\theta : r(\theta)\le 0\}$ (e.g.,
box constraints, effect-ordering, sign/monotonicity) replace the unbounded
behavior in the null space of Proposition~\ref{prop:dichotomy} with a finite LP
whose width we can compute and optimize.

\subsection{A tractable surrogate}
\label{sec:surrogate}

Fix a candidate $k$ and a target $(t_m, t_r)$. The width is the value of the LP
pair \eqref{eq:bounds} with $\Sigma \leftarrow \Sigma_{\mathcal{S}\cup\{k\}}$.
Because $\Delta\Sigma_k$ has rank at most $2$, the candidate's effect on the
feasible set is a \emph{rank-2 perturbation} of the affine subspace
$\{\Sigma_\mathcal{S}\theta = -b\}$. Under Assumption~6 of \citet{boyarsky2026} (Robinson constraint qualification) the LP values are Lipschitz in
$\Sigma$, and we obtain a first-order surrogate:

\begin{proposition}[Rank-2 envelope]
\label{prop:envelope}
Let $(\theta_\mathcal{S}^\star, \lambda_\mathcal{S}^\star,
\nu_\mathcal{S}^\star)$ be a KKT triple at $\Sigma_\mathcal{S}$ for a target
$(t_m, t_r)$, with minimum-norm multiplier $\lambda_\mathcal{S}^\star$. Then
\begin{equation}
\label{eq:envelope}
v_u\bigl(\gamma_{\mathcal{S}\cup\{k\}}\bigr) - v_u\bigl(\gamma_\mathcal{S}\bigr)
 \;\approx\;
 - \lambda_\mathcal{S}^{\star\top}\bigl(\Delta\Sigma_k\,
 \theta_{\mathcal{S}}^\star + \Delta b_k\bigr),
\end{equation}
and analogously for $v_\ell$. The first-order change in the width is
$\Delta w_k \approx - (\lambda_\mathcal{S}^{u,\star} - \lambda_\mathcal{S}^{\ell,\star})^\top
 (\Delta\Sigma_k\,\theta^\star_{\mathcal{S}} + \Delta b_k)$.
\end{proposition}

Proposition~\ref{prop:envelope} is the Bonnans--Shapiro directional derivative of
the constrained-value map, specialized to the finite-dimensional rank-2
site-update direction $(h_\Sigma, h_b) = (\Delta\Sigma_k/\|\Delta\Sigma_k\|,
\Delta b_k)$. The formula \eqref{eq:envelope} coincides with Equation~(1.3) of \citet{boyarsky2026}, evaluated at the site-augmentation direction. It gives a
\emph{linear-in-$k$} surrogate for the width change, so the minimax selection
\eqref{eq:minimax} reduces to a finite computation:
\begin{equation}
\label{eq:surrogate}
k^\star_{\rm surr}\;\in\;\arg\min_{k\in\mathcal{K}}\;
 \sup_{(t_m,t_r)\in\mathcal{T}}\;
 \Bigl|\,\lambda_\mathcal{S}^{u,\star}(t)^\top
  \Delta\Sigma_k\,\theta^\star_\mathcal{S}(t)
  - \lambda_\mathcal{S}^{\ell,\star}(t)^\top
  \Delta\Sigma_k\,\theta^\star_\mathcal{S}(t)\,\Bigr|,
\end{equation}
after taking a sup over the unknown $\Delta b_k$ in the prior ball $\mathcal{B}$,
which is handled by a standard dual-norm bound and is additive in the candidate.
Computing \eqref{eq:surrogate} requires solving the LP pair \eqref{eq:bounds} at
$\Sigma_\mathcal{S}$ once per target and a finite-dimensional matrix-vector
product per candidate, total cost $O(|\mathcal{T}|\cdot K)$.

\subsection{Three regimes}
\label{sec:regimes}

The linearization \eqref{eq:envelope} inherits the three-regime taxonomy of the
LOO diagnostic of \citet[\S4.4]{boyarsky2026}, now read in reverse:

\begin{itemize}
\item \textbf{(A) Coverage-expanding candidate.} If $(C_m^{(k)}, C_r^{(k)})
  \notin \mathrm{aff}(\mathcal{H}_\mathcal{S})$, then
  $\mathrm{null}(\Sigma_{\mathcal{S}\cup\{k\}}) \subsetneq
  \mathrm{null}(\Sigma_\mathcal{S})$. Proposition~\ref{prop:envelope} detects
  this as a first-order contraction of the width whenever the binding recession
  direction at target $t$ lies along the newly-resolved null-space direction. In
  the unrestricted limit this regime collapses infinite widths to finite
  ones.

\item \textbf{(B) Location-improving candidate.} If $(C_m^{(k)}, C_r^{(k)}) \in
  \mathrm{aff}(\mathcal{H}_\mathcal{S})$ but outside the convex hull, the rank of
  $\Sigma$ does not change but $\Delta b_k$ translates the moment-implied affine
  subspace. The binding vertex of the LP can shift, which gives a first-order
  width change of either sign. This is the regime in which a new site
  improves
  overlap at the target without adding a new identification direction.

\item \textbf{(C) Feasibility-restoring (or -breaking) candidate.} If $\{\theta
  : \Sigma_\mathcal{S}\theta = -b_\mathcal{S}\}\cap \mathcal{R} = \emptyset$, an
  infeasibility signaled by the Farkas alternative of \citet[\S4.4 (c)]{boyarsky2026}, a candidate $k$ can restore feasibility if its site-update
  direction pushes the affine subspace back into $\mathcal{R}$.
  Conversely, a candidate whose direction is incompatible with the existing
  restriction set can break feasibility, which is diagnostic information rather
  than a reason to exclude the candidate.
\end{itemize}

The minimax selection \eqref{eq:minimax} therefore does not reduce to a single
geometric quantity in general. The convex-hull rule is exact in regime (A) and
the unrestricted limit, and it is a useful first screen that should be refined
by \eqref{eq:surrogate} when the restriction set is binding.

\subsection{Handling the unknown \texorpdfstring{$\Delta b_k$}{Delta b\_k}}
\label{sec:unknown-b}

The right-hand-side update $\Delta b_k$ depends on the candidate site's future
outcomes. Three prior choices are natural:

\begin{enumerate}
\item \textbf{Uninformative ball.} $\Delta b_k \in \{b : \|b\|_2 \le B\}$. Yields
  a worst-case width via dual-norm bound, but can dominate $\Delta\Sigma_k$ at
  small candidate sample sizes. Conservative.

\item \textbf{Meta-analytic prior.} Use the existing-pool estimates as a prior on
  the candidate's $(Y, A)$ conditional means, and propagate through the
  expression for $b$ in Equation~(A.7) of \citet{boyarsky2026}. This is a
  structural Bayesian design; under a flat prior on $\theta$ it recovers a
  Bayesian D-optimality criterion, and the minimax and Bayesian selections agree
  in the large-candidate-sample limit.

\item \textbf{Adversarial $b$.} $\Delta b_k \in \mathcal{B}$ chosen to maximize
  the width. This is the robust-optimization analogue of (1) with an ellipsoidal
  $\mathcal{B}$; Sion-minimax arguments give a tractable saddle-point
  representation.
\end{enumerate}

In our simulation (Section~\ref{sec:simulation}) we adopt (1) with $B$ set from
the existing-pool residual variance. The three variants agree on the
minimax-optimal candidate in all our runs. This is expected, because the rank
update in \eqref{eq:dsigma} dominates the first-order width change in the
regimes where candidates differ meaningfully.
